\chapter{Implementação do Algoritmo XMSS em Python}

Clonei o repositório do liboqs. Compilei a biblioteca como Shared Library.
Para implementar o algoritmo XMSS em Python utilizando a liboqs, é necessário fazer a ponte entre a biblioteca compilada em C e o ambiente Python, além de utilizar a API específica para esquemas stateful (\textit{sig\_stfl.h}. Clonei o repositório liboqs-python, que é o wrapper da biblioteca para a linguagem python. Vincular a biblioteca compartilhada ao SO executando o comando:

% export LD_LIBRARY_PATH=~/Documents/liboqs/build/lib:$LD_LIBRARY_PATH
export LD\_LIBRARY\_PATH=~/Documents/liboqs/build/lib:\$LD\_LIBRARY\_PATH

Navegar até o repostiório do liboqs-python, Criar e ativar o ambiente virtual:

% python -m venv .venv
python -m venv .venv

%source .venv/bin/activate
source .venv/bin/activate

Instalar o wrapper:

%pip install .
pip install .


Executei um teste para ver se a biblioteca era reconhecida mas o liboqs-python foi projetado para suportar apenas as APIs de algoritmos Stateless, vou poder usar na implementação do SPHINCS+, por agora terei que usar ctypes.
:q
:q
